\chapter{系统功能实现}

本章首先介绍前后端工程目录结构，随后依照功能模块对平台的实现情况展开说明，并结合关键代码片段、业务流程图以及系统运行截图，把各模块的核心功能与实现细节进行展示。

\section{工程目录结构}

清晰的工程目录结构是保证代码可维护性的重要前提。本系统采用前后端分离架构，两端目录按职责独立组织。

\subsection{前端工程结构}

前端基于 Vite 脚手架初始化，源码统一置于 \texttt{frontend/src/} 目录下，按如下方式分层组织：

\begin{lstlisting}[numbers=none]
frontend/src/
├── api/            # Axios 请求封装，按模块拆文件
│   ├── auth.js     # 认证相关接口（登录/注册）
│   ├── student.js  # 学生端接口
│   ├── counselor.js# 咨询师端接口
│   └── school.js   # 管理端接口
├── layouts/        # 三端公共布局组件
├── router/         # Vue Router 路由配置与路由守卫
├── stores/         # Pinia 全局状态（登录信息/Token）
├── utils/          # 工具函数（统一请求拦截 request.js）
└── views/
    ├── student/    # 学生端页面（6 个视图）
    ├── counselor/  # 咨询师端页面（5 个视图）
    └── school/     # 管理端页面（3 个视图）
\end{lstlisting}

\subsection{后端工程结构}

后端遵循 Controller $\rightarrow$ Service $\rightarrow$ Model 三层分离原则，源码置于 \texttt{backend/src/} 目录下：

\begin{lstlisting}[numbers=none]
backend/src/
├── config/         # 数据库连接配置
├── controllers/    # 请求处理层（解析请求、组装响应）
│   ├── auth.controller.js
│   ├── student.controller.js
│   ├── counselor.controller.js
│   ├── school.controller.js
│   └── ai.controller.js
├── middleware/     # JWT 验证（auth.js）、角色守卫（roleGuard.js）
├── models/         # Sequelize 模型定义（10 张数据表）
├── routes/         # 模块化路由，按角色拆分文件
├── services/       # 业务逻辑层（含 ai.service.js）
└── utils/          # JWT 工具函数、统一响应格式等
\end{lstlisting}

这种分层目录使每个模块职责单一：Controller 只负责请求解析与响应组装，业务逻辑下沉到 Service，模型层只描述数据结构，当某一层需要修改时不会波及其他层。

\section{用户认证与权限模块}

认证模块是整个平台开展安全控制的基础部分，负责学生、咨询师以及管理员三类角色的注册、登录以及全过程鉴权工作，并借助 JSON Web Token（JWT）来实现无状态认证；所有受保护接口都需要经过认证中间件以及角色守卫的双重校验，这样一来可以确保不同角色之间的数据访问得到安全隔离。

当用户在登录页面输入账号和密码后，系统会先对其身份开展验证工作，在校验通过之后签发 JWT，前端则把 Token 写入本地存储，后续发起的各类请求都会自动携带这一凭证来完成鉴权。图~\ref{fig:4-3}~给出了平台的登录界面，页面顶部设置了角色切换按钮，支持学生、咨询师以及管理员三端分别进行登录；登录完成后，系统会按照角色自动跳转到对应首页，并且鉴于三端路由彼此隔离，任何越权访问都会被重定向回本角色首页。

\textbf{后端认证中间件。}代码~\ref{lst:auth-mid}~给出了 JWT 认证中间件的核心实现。中间件从请求头中提取 Bearer Token，调用 \texttt{verifyToken} 完成签名校验，校验通过后将用户 \texttt{id} 与 \texttt{role} 挂载到 \texttt{req.user}，供后续中间件及 Controller 直接使用。

\begin{lstlisting}[language=JavaScript, caption={JWT 认证中间件（backend/src/middleware/auth.js）}, label={lst:auth-mid}]
function authMiddleware(req, res, next) {
  const authHeader = req.headers['authorization'];
  if (!authHeader || !authHeader.startsWith('Bearer '))
    return res.status(401).json({ code: 401, message: '(*@未提供认证 Token@*)' });
  try {
    const payload = verifyToken(authHeader.slice(7));
    req.user = payload; // (*@挂载用户 id 与 role@*)
    next();
  } catch (err) {
    const msg = err.name === 'TokenExpiredError'
      ? '(*@Token 已过期@*)' : '(*@Token 无效@*)';
    return res.status(401).json({ code: 401, message: msg });
  }
}
\end{lstlisting}

\textbf{角色守卫中间件。}在认证中间件校验身份之后，\texttt{roleGuard} 进一步校验角色，代码如下。路由定义时可灵活组合，例如 \texttt{roleGuard('student')} 仅允许学生访问，\texttt{roleGuard('counselor', 'school')} 则同时放行咨询师和管理员。

\begin{lstlisting}[language=JavaScript, caption={角色守卫中间件（backend/src/middleware/roleGuard.js）}, label={lst:role-guard}]
function roleGuard(...roles) {
  return (req, res, next) => {
    if (!roles.includes(req.user.role))
      return res.status(403).json({ code: 403, message: '(*@无权限访问@*)' });
    next();
  };
}
\end{lstlisting}

\textbf{前端 Axios 请求拦截。}代码~\ref{lst:axios}~给出了前端统一请求封装的核心逻辑。请求拦截器在每次发出请求前自动从 \texttt{localStorage} 读取 Token 并写入请求头；响应拦截器在收到 401 状态码时自动清除本地凭据并跳转登录页，实现了认证失效的统一处理。

\begin{lstlisting}[language=JavaScript, caption={Axios 请求拦截器（frontend/src/utils/request.js）}, label={lst:axios}]
request.interceptors.request.use(config => {
  const token = localStorage.getItem('token');
  if (token) config.headers.Authorization = `Bearer ${token}`;
  return config;
});
request.interceptors.response.use(
  res => res.data,
  error => {
    if (error.response?.status === 401) {
      localStorage.removeItem('token');
      router.push('/login');
    }
    return Promise.reject(new Error(error.response?.data?.message));
  }
);
\end{lstlisting}

\begin{figure}[H]
    \centering
    \includegraphics[width=0.85\textwidth]{fig4-3.png}
    \caption{用户登录界面}
    \label{fig:4-3}
\end{figure}

\section{学生端功能模块}

学生端是平台当中的核心服务入口，围绕"发现咨询师 $\rightarrow$ 发起预约 $\rightarrow$ AI 辅助对话 $\rightarrow$ 评价反馈"这一整套使用路径，持续为学生提供覆盖全过程的心理健康服务支持。

\subsection{咨询师浏览与预约}

学生在咨询师列表页面可以查看全部处于可预约状态的咨询师姓名、职称、专业方向、擅长问题类型以及可预约时段，并在此基础上完成目标咨询师的筛选以及预约发起。图~\ref{fig:4-4}~给出了该列表页面，页面中的每张卡片都会较为直观地展示咨询师的核心信息，学生可直接点击"预约"按钮进入对应的预约页面。

\begin{figure}[H]
    \centering
    \includegraphics[width=0.85\textwidth]{fig4-4.png}
    \caption{学生端咨询师浏览与预约界面}
    \label{fig:4-4}
\end{figure}

预约提交后进入"待确认"状态，等待咨询师审核。在预约历史页，学生可查看历次预约的完整状态流转记录（待确认、已确认、已完成、已取消），对尚未开始的预约可主动取消。咨询会话完成后，学生可提交 1—5 分评分及文字反馈，评价数据汇入管理端统计系统。

\subsection{公告查阅}

学生端首页会集中展示由学校管理端所发布的最新公告，这些公告会依照发布时间倒序进行排列，内容包括心理健康活动通知、咨询中心动态以及学校政策说明，以便帮助学生及时了解校园心理服务信息。

在数据层面，公告内容存储于 \texttt{announcements} 表，学生端通过 \texttt{GET /api/v1/student/announcements} 接口获取全量公告列表，接口返回按 \texttt{created\_at} 降序排列的结果，无需额外分页参数即可完成首屏加载。咨询师端同样可以查阅公告，确保一线服务人员与学生获取相同的通知信息，实现全平台信息同步。公告内容由学校管理端统一维护，管理员在发布时填写标题与正文，发布即时生效；对于过期公告，管理员可执行删除操作，系统在删除前弹出二次确认提示，防止误操作。这一设计使校园心理服务通知能够以标准化方式触达师生，无需借助其他即时通讯渠道进行补充传达。

\section{咨询师端功能模块}

咨询师端以"预约审核 $\rightarrow$ 开展会话 $\rightarrow$ 记录归档"这一核心工作流程为主线，面向咨询师提供案例管理方面的功能支持，并且使相关操作数据能够与学生端实现实时联动。

\subsection{预约审核与会话管理}

图~\ref{fig:4-2}~对预约申请直至会话完成的完整业务流程进行了展示。咨询师会先在工作台中查看近期提交的预约申请，并在逐项对来访者信息进行审阅之后，作出确认或者拒绝的处理决定。如果预约得到确认，对应状态就会转为"已确认"；待约定时间到达后，咨询师创建会话记录并正式开展会谈。会谈结束之后，咨询师执行"结束会话"操作，系统还会同步把预约状态更新为"已完成"，由此形成完整的咨询记录闭环。

\begin{figure}[H]
    \centering
    \includegraphics[width=0.85\textwidth]{fig4-2.png}
    \caption{预约与会话业务流程图}
    \label{fig:4-2}
\end{figure}

在具体实现这一层面上，预约状态每一次发生流转时，都会由后端接口来开展统一控制，并且借助数据库事务保证相关操作的原子性。以"结束会话"这一场景为例，后端会在同一个事务当中同步把 \texttt{sessions.status} 置为 \texttt{completed}，同时把 \texttt{appointments.status} 置为 \texttt{completed}；如果其中任一步骤执行失败，系统就会进行整体回滚，这样一来能够保证两张表的状态始终保持一致，避免出现"会话已结束但预约仍显示进行中"的数据不一致情况。

\subsection{案例笔记与工作台}

围绕每一次会话，咨询师都可以去撰写私密的案例笔记，用来记录来访者的主诉要点、情绪状态评估、会谈结论以及后续干预计划；该笔记支持进行多次追加，并且只对咨询师本人可见。图~\ref{fig:4-6}~给出了咨询师工作台首页，页面借助统计卡片以及快速操作入口，把待处理任务进行集中展示，以便咨询师及时把握整体工作情况。

\begin{figure}[H]
    \centering
    \includegraphics[width=0.85\textwidth]{fig4-6.png}
    \caption{咨询师端工作台界面}
    \label{fig:4-6}
\end{figure}

\section{学校管理端功能模块}

学校管理端主要面向心理中心管理人员来使用，拥有咨询师账号管理、平台数据统计以及公告发布这三项运营管理功能，同时不会直接介入具体咨询内容，从而对咨访关系的私密性加以保障。

\subsection{咨询师账号管理}

管理员可以对咨询师账号进行创建，并完成初始密码的设置工作；在账号创建完成之后，咨询师即可登录系统，对个人档案信息进行维护。对于已经离职或者暂时停止服务的咨询师，系统能够借助禁用处理使其不再于学生端显示，同时保留历史会话数据的完整记录，便于后续开展查阅。

在实现层面，账号创建接口接收用户名、邮箱和初始密码，后端在写库之前对密码统一进行 bcrypt 哈希处理，并将 \texttt{role} 字段设置为 \texttt{counselor}，\texttt{is\_active} 默认为 \texttt{1}（启用）。禁用操作调用独立的 \texttt{PATCH} 接口，仅将目标账号的 \texttt{is\_active} 字段更新为 \texttt{0}，而不删除任何相关数据行；学生端在获取可预约咨询师列表时，查询条件中会附加 \texttt{is\_active = 1} 的过滤，因此被禁用的咨询师将自动从列表中消失，同时其历史预约与会话记录仍完整保留，管理员可随时重新启用。这种"软禁用"策略在兼顾数据完整性的同时，也保留了对咨询师账号进行灵活调整的能力。

\subsection{平台数据统计与公告管理}

数据统计页面分为两个层次。第一层是平台概览卡片，集中展示在读学生人数、在岗咨询师人数、累计预约总量、已完成会话数量以及平均满意度等核心指标，以卡片形式直观呈现。第二层是图表统计区，借助 Apache ECharts 库渲染三张可视化图表：其一为各咨询师预约量环形饼图，后端对 \texttt{appointments} 表按 \texttt{counselor\_id} 分组聚合，并关联咨询师真实姓名后返回；其二为学生预约量 Top 10 柱状图，按预约次数降序取前十名，以学生姓名与学号联合标注横轴；其三为咨询话题分布环形饼图，统计各话题类别（学业压力、情感问题、人际关系、家庭问题、就业焦虑、心理健康、其他）的预约占比。为支撑话题维度的统计，\texttt{appointments} 表新增了枚举类型的 \texttt{topic} 字段，学生在发起预约时通过下拉选择器指定话题类别，数据落库后即可参与聚合统计。图~\ref{fig:4-7}~给出了管理端的数据概览界面，相关指标均取自系统的实际运行数据。

\begin{figure}[H]
    \centering
    \includegraphics[width=0.85\textwidth]{fig4-7.png}
    \caption{学校管理端数据概览界面}
    \label{fig:4-7}
\end{figure}

在公告管理这一环节中，管理员在完成标题和正文内容的填写之后，便可以将公告发布出去，公告一经发布，就会马上在学生端以及咨询师端同步展示；对于已经不再适用的历史公告，管理员也可以进行删除操作，并且系统会在删除之前借助弹框开展二次确认，从而避免出现误操作。

\section{AI 心理助手模块}

AI 心理助手属于平台当中的特色功能模块，以"心语"作为名称，并且依托 DeepSeek 大语言模型来进行构建，具有情感倾听、心理疏导以及咨询知识解答等方面的能力，能够持续为学生提供 7×24 小时的情绪支持服务。对于涉及自伤、自杀倾向等严重心理危机的情况，助手需要主动把学生引导至专业咨询师或者心理援助热线，由此把 AI 辅助和专业心理治疗之间的边界加以明确。

\subsection{流式对话交互效果}

系统依靠 Server-Sent Events（SSE）技术来实现 AI 回复内容的流式输出；在用户发送消息之后，AI 所生成的回复会以"打字机效果"的形式逐词显示，因此无须等到完整回复全部生成，便可以实时查看输出内容，这样一来可明显减轻交互过程中的等待感。图~\ref{fig:4-1}~展示了 AI 流式对话的完整交互时序。

\begin{figure}[H]
    \centering
    \includegraphics[width=0.85\textwidth]{fig4-1.png}
    \caption{AI 流式对话交互时序图}
    \label{fig:4-1}
\end{figure}

{\sloppy
在具体的技术实现层面，服务端依靠设置响应头 \texttt{Content-Type:~text/event-stream} 来建立 SSE 通道，并调用 \texttt{res.flushHeaders()} 把响应头马上发送到客户端；在这之后，每当 DeepSeek API 返回一个文字片段，后端就会调用 \texttt{res.write()}，把它封装成标准的 SSE 事件格式后推送给前端，所以整个过程不需要等完整回复生成。前端则以 \texttt{ReadableStream} 方式去读取响应体，每接收到一个数据块，都会触发 Vue 的响应式更新，把新字符追加到消息气泡当中，进而形成文字逐词显现的视觉效果；当流式传输结束以后，服务端会推送 \texttt{data:~[DONE]} 信号，客户端随即关闭连接，并把完整回复内容持久化到本地状态。
}

代码~\ref{lst:sse}~给出了服务端建立 SSE 通道并实时转发流式片段的核心实现（摘自 \texttt{services/ai.service.js}）。

\begin{lstlisting}[language=JavaScript, caption={SSE 流式通道建立与片段转发}, label={lst:sse}]
// (*@设置 SSE 响应头，建立持久流式通道@*)
res.setHeader('Content-Type', 'text/event-stream');
res.setHeader('Cache-Control', 'no-cache');
res.setHeader('Connection', 'keep-alive');
res.flushHeaders();

// (*@逐块读取 DeepSeek 流式响应并即时推送前端@*)
const { done, value } = await reader.read();
const delta = parsed.choices?.[0]?.delta?.content;
if (delta) {
  res.write('data: ' + JSON.stringify({ type: 'delta', content: delta }) + '\n\n');
}
// (*@流式传输结束后发送完成信号@*)
res.write('data: ' + JSON.stringify({ type: 'done' }) + '\n\n');
res.end();
\end{lstlisting}

\subsection{多轮对话上下文管理}

系统会自动对最近 20 条对话消息进行维护，并把它们作为上下文来使用，从而使 AI 在多轮交流当中能够理解前面的内容并保持回应的连贯性。图~\ref{fig:4-5}~展示了学生与 AI 心理助手开展对话的真实运行效果，AI 可以结合学生描述的具体情境，比如学业压力以及焦虑情绪，给出有针对性的舒缓建议；历史对话列表会按照最近活跃时间进行排序，便于随时查看与回顾。

\begin{figure}[H]
    \centering
    \includegraphics[width=0.85\textwidth]{fig4-5.png}
    \caption{AI 心理助手对话界面}
    \label{fig:4-5}
\end{figure}

在上下文管理的实现过程当中，系统每次接收到用户发送的新消息之后，会先把该消息写入 \texttt{ai\_messages} 表，随后依靠 \texttt{ORDER BY created\_at ASC LIMIT 20} 检索当前会话最近 20 条历史消息，并按照时间先后顺序组装成 \texttt{messages} 数组，再提交给 DeepSeek API。如果对话轮次已经超过 20 条，那么最早的一条消息就会自动移出上下文窗口，但仍然保存在数据库当中，以便用户后续对历史记录进行查阅。这种"固定滑动窗口"方式在保持对话连贯性的同时，还把每次 API 调用的 Token 消耗控制在较为合理的范围内，从而较好平衡了对话质量以及推理成本。

\section*{本章小结}

本章首先对前后端工程目录结构进行了说明，从整体上展示了系统按职责分层组织代码的方式；随后围绕用户认证、学生端、咨询师端、学校管理端以及 AI 心理助手这五个功能模块，结合关键代码片段、业务流程图以及界面截图，把平台核心功能的实现细节进行了较为系统的展示。认证模块借助 JWT 认证中间件与角色守卫中间件，完成了三端数据访问方面的安全隔离，并通过 Axios 拦截器在前端统一管理 Token 生命周期；学生端按照"浏览—预约—对话—评价"的完整链路，推动心理咨询服务全流程得以数字化；咨询师端构建了一个覆盖审核到归档环节的一体化案例管理工具；管理端拥有账号管理（软禁用策略）、数据统计以及公告发布等学校运营支持功能；AI 心理助手则凭借 SSE 流式对话技术，把大语言模型能力接入平台，进而为学生提供随时可获得的情绪支持入口。
